\section{Evaluation}
At the conclusion of the iTelos procedure we obtained a
Knowledge Graph (KG) generated from external data of the real world;
however it is important now to proceed by understanding
if the result satisfies our needs and expected performances,
for this it is necessary to make an evaluation of the KG,
and in order to do it we have to rely on a software that
interrogates our KG: such tool
is represented by \href{https://www.ontotext.com/products/graphdb/}{GraphDB}, a program that allows
to load the final KG and execute SPARQL queries on it in order
to retrieve the necessary informations. In this final
section we will cover three different aspects of the 
evaulation process:
\begin{itemize}
	\item the general statistics of our KG
	\item the knowledge evaluation phase,
		in particular the different coverage types
	\item the data evaluation phase, foucsing
		mainly on the connectivity of the KG
	\item the conversion of the various
		Competency Questions in SPARQL queries,
		in order to check if our KG is able to
		satisfy them
\end{itemize}

\subsection{KG statistics}

As first point in the evaluation it is useful to report some
statistics in order to have a better understanding of the situation:

\begin{itemize}
	\item In total there are 12 etypes, and among them we count: \textit{MeteorologicalMeasurement},
		\textit{TempMeasurement},\textit{Month},\textit{WindMeasurement},
		\textit{RainMeasurement},\textit{HumidityMeasurement},\textit{PressureMeasurement},
		\textit{SolarRadiationMeasurement},\textit{MeteorologicalFacility},\textit{Valley},
		\textit{Coordinates},\textit{Sensor}
	\item 12 data properties and 10 object properties		
\end{itemize}

Both these statistics have been retrieved through following SPARQL queries
that can be found \href{https://github.com/zerise000/KG_Project2025_ClimateChange/tree/main/Evaluation/KG_statistics}{here}

\subsection{Knowledge Layer evaluation}

An important aspect of Knowledge Graphs is how well they are able
to cover the portion of knowledge we want to incapsulate inside
them, specifically this concept can be illustrated with the Competency
Questions (CQs): if a particular Knowledge Graph has been designed well
then it will be able to cover as many CQs as possible. We formalize
this concept by saying that we want to compute the \textbf{coverage}
of our teleontology $T$ at Knowledge Level. Referencing the previous
sections we can retrieve:

\begin{itemize}
	 
	\item  the \textbf{Entity Types Coverage}, which measures the ratio
	of how many etypes are present in our teleontology: 
		$$Cov_E(CQ_E) = \frac{|CQ_E \cap T_E|}{|CQ_E|} = \frac{11}{18} = 61.1\%$$
	where $T_e$ are the etypes in our generated teleontology, while $CQ_e$
	are the etypes related to the Competency Questions $CQs$.

	\item last, the \textbf{Property Coverage} it measures
	how many attributes in the Knowledge Layer are present in the
	Teleology. The formula is:
		$$Cov_{P}(CQ_{P}) = \frac{|CQ_{P} \cap T_{P}|}{|CQ_{P}|} = \frac{12}{29} = 41.3\%$$
	where $T_{P}$ are the properties in our generated teleontology, while $CQ_{P}$
	are the properties associated to the Comptetency Questions $CQs$.
\end{itemize}

It is also possible to define the coverage by analyzing the teleontology
$T$ generated in the previous sections, and comparing it with various 
knowledge elements called Knowledge Atoms (KATOMs) taken from
the reference ontology generated in the knowledge phase:

\begin{itemize}
	\item when the considered KATOM is an etype we talk about
	\textbf{Entity Type Coverage}, which measures the ratio
	of how many etypes are present in our teleontology: 
		$$Cov_e(KATOM_e) = \frac{|KATOM_e \cap T_e|}{|KATOM_e|} = \frac{12}{18} = 66.6\%$$
	where $T_e$ are the etypes in our generated teleontology, while $KATOM_e$
	are the etypes defined at knowledge level.

	\item similarly, we define the \textbf{Object Property Coverage} 
	when the considered KATOM is about relationships, and it measures
	how many relationships in the Knowledge Layer are present in the
	Teleology. In this case the formula becomes:
		$$Cov_{OP}(KATOM_{OP}) = \frac{|KATOM_{OP} \cap T_{OP}|}{|KATOM_{OP}|} = \frac{5}{12} = 41.6\%$$
	where $T_{OP}$ are the relationships in our generated teleontology, while $KATOM_{OP}$
	are the relationships defined at knowledge level.

	\item last, the \textbf{Data Property Coverage} treats 
	KATOMs as attributes, and it measures
	how many attributes in the Knowledge Layer are present in the
	Teleology. The formula is:
		$$Cov_{DP}(KATOM_{DP}) = \frac{|KATOM_{DP} \cap T_{DP}|}{|KATOM_{DP}|} = \frac{11}{13} = 84.61\%$$
	where $T_{DP}$ are the relationships in our generated teleontology, while $KATOM_{DP}$
	are the relationships defined at knowledge level.
\end{itemize}

All the SPARQL queries used to obtain the various
Teleontology vs KATOM coverages can be found here.
As we can see the Object Property Coverage results to be low, which means that there is not much
correlation between the different etypes of the Knowledge Graph. One of the possible causes
could be the usage of a single informal dataset for the building phase instead on relying on
multiple sources. The Entity Type coverage also seems to have some gaps inside it, in fact
we can identify the missing etypes \href{https://github.com/zerise000/KG_Project2025_ClimateChange/blob/main/Evaluation/knowledgeEvaluation/missingEtypes.txt}{through the following SPARQL query}:

\begin{itemize}
	\item ClimateChange
	\item Trentino
	\item TimeInterval
	\item ClimateChangeIndicator
	\item Trend
\end{itemize}

\subsection{Data Layer evaluation}

the knowledge layer evaluation gives us only a generic view on the situation:
in fact it tells us how much knowledge is covered by the defined KG, however
it does not provide any useful informations about the structure of the graph
itself, like the connectivity between the nodes. The evaluation of the data
layer aims to fill this gap by representing the 
Entity and Property connectivity through a \textbf{Connectivity Matrix}.
\begin{table}
	\centering
	\resizebox{\textwidth}{!}{%
	\begin{tabular}{| l | l | l | l | l | l | l | l | l | l | l | l | l |}
		\hline
		Etypes & \textbf{MeteorologicalMeasurement} & \textbf{TempMeasurement} & \textbf{Month} & \textbf{WindMeasurement} &
		\textbf{RainMeasurement} & \textbf{HumidityMeasurement} & \textbf{PressureMeasurement} & \textbf{SolarRadiationMeasurement} & 
		\textbf{MeteorologicalFacility} & \textbf{Valley} & \textbf{Coordinates} & \textbf{Sensor} \\
		\hline
		\textbf{MeteorologicalMeasurement} & 36101 & 0 & 14541 & 0 & 0 & 0 & 0 & 0 & 0 & 0 & 0 & 0 \\ 
		\hline
		\textbf{TempMeasurement} & 0 & 17610 & 6106 & 0 & 0 & 0 & 0 & 0 & 0 & 0 & 0 & 0 \\
		\hline
		\textbf{Month} & 0 & 0 & 2806 & 0 & 0 & 0 & 0 & 0 & 0 & 0 & 0 & 0 \\
		\hline
		\textbf{WindMeasurement} & 0 & 0 & 2572 & 7240 & 0 & 0 & 0 & 0 & 0 & 0 & 0 & 0 \\
		\hline
		\textbf{RainMeasurement} & 0 & 0 & 2447 & 0 & 4776 & 0 & 0 & 0 & 0 & 0 & 0 & 0 \\
		\hline
		\textbf{HumidityMeasurement} & 0 & 0 & 1183 & 0 & 0 & 2247 & 0 & 0 & 0 & 0 & 0 & 0 \\
		\hline
		\textbf{PressureMeasurement} & 0 & 0 & 1141 & 0 & 0 & 0 & 2163 & 0 & 0 & 0 & 0 & 0 \\
		\hline
		\textbf{SolarRadiationMeasurement} & 0 & 0 & 1092 & 0 & 0 & 0 & 0 & 2065 & 0 & 0 & 0 & 0 \\
		\hline
		\textbf{MeterelogicalFacility} & 0 & 0 & 0 & 0 & 0 & 0 & 0 & 0 & 36101 & 0 & 24 & 23 \\
		\hline
		\textbf{Valley} & 0 & 0 & 0 & 0 & 0 & 0 & 0 & 0 & 0 & 24 & 0 & 0 \\
		\hline
		\textbf{Coordinates} & 0 & 0 & 0 & 0 & 0 & 0 & 0 & 0 & 0 & 24 & 72 & 0 \\
		\hline
		\textbf{Sensor} & 14541 & 6106 & 0 & 2572 & 2447 & 1183 & 1141 & 1092  & 0 & 0 & 0 & 0 \\
		\hline
		\end{tabular}%
		}
\end{table}
In particular we define two metrics:

\begin{itemize}
	\item the \textbf{Entity Connectivity} for a specific EType X, which
	represents how much entities are connected to each other: $EC(X) = \frac{\sum^N_{Y=1}{(X,Y)}}{OP(X)}$
	\item the \textbf{Property Connectivity} for a specific EType X, which tells us
	how much the entities are connected to their properties: $PC(X) = \frac{(X,X)}{DP(X)}$
\end{itemize}

Where $(X,Y)$ is a cell in the connectivity matrix, $OP(X)$ is the number of object properties
for the EType $X$, and $DP(X)$ is the number of data properties in the EType $X$. By calculating
the number of object properties and data properties for each etype we can easily retrieve 
the Entity Connectivitie:
\begin{itemize}
	\item \textit{MeteorologicalMeasurement}: $\frac{36101+14541}{1} = 50642$
	\item \textit{TempMeasurement}: $\frac{17610+6106}{1} = 23716$
	\item \textit{Month}: $0$
	\item \textit{WindMeasurement}: $\frac{2572+7240}{1} = 9812$
	\item \textit{RainMeasurement}: $\frac{4776+2447}{1} = 7223$
	\item \textit{HumidityMeasurement}: $\frac{2247+1183}{1} = 3430$
	\item \textit{PressureMeasurement}: $\frac{1141+2163}{1} = 3304$
	\item \textit{SolarRadiationMeasurement}: $\frac{2065+1092}{1} = 3157$
	\item \textit{MeteorologicalFacility}: $\frac{36101+24+23}{2} = 18074$
	\item \textit{Valley}: $0$
	\item \textit{Coordinates}: $\frac{72+24}{1} = 96$
	\item \textit{Sensor}: $\frac{14541+6106+2572+2447+1183+1141+1092}{1} = 29082$
\end{itemize}
Also we can find the Property Connectivity of each etype:
\begin{itemize}
	\item \textit{MeteorologicalMeasurement}: $\frac{36101}{4} = 9025$
	\item \textit{TempMeasurement}: $\frac{17610}{4} = 4402$
	\item \textit{Month}: $\frac{2806}{1} = 2806$
	\item \textit{WindMeasurement}: $\frac{7240}{3} = 2413$
	\item \textit{RainMeasurement}: $\frac{4776}{2} = 2388$
	\item \textit{HumidityMeasurement}: $\frac{2247}{2} = 1123$
	\item \textit{PressureMeasurement}: $\frac{2163}{2} = 1081$
	\item \textit{SolarRadiationMeasurement}: $\frac{2065}{2} = 1032$
	\item \textit{MeteorologicalFacility}: $\frac{36101}{4} = 9025$
	\item \textit{Valley}: $\frac{24}{1} = 24$
	\item \textit{Coordinates}: $\frac{72}{3} = 24$
	\item \textit{Sensor}: $0$
\end{itemize}
We can notice how the 
Property Connectivity of Sensor is set to 
zero since there is no instance of this etype 
in our Knowledge Graph.
Both these metrices can be generalized to the whole knowledge graph:

\begin{itemize}
	\item the \textbf{Entity Connectivity}: $EC(X) = \sum^{N}_{X=1}{EC(X)} = 148536$
	\item the \textbf{Property Connectivity}: $PC(X) = \sum^{N}_{X=1}{PC(X)} = 33343$
\end{itemize}

\subsection{Queries}

In previous sections we defined Competency Questions (CQs) that describes hypotetical scenario
that our Knowledge Graph (KG) should be able to cover, which must be then turned into SPARQL queries
and run on our teleontology in order to check it returns the predicted results. We can model
the whole situation in the following way: for each CQ we describe
values it should display and a description of the actual results:
\begin{itemize}
	
	\item The Competency Question Q1 query can be found \href{https://github.com/zerise000/KG_Project2025_ClimateChange/blob/main/Evaluation/competencyQuestions/Q1.txt}{here}. 
	From the theoretical point of view of our model the query should
	return the temperatures of different metereological stations in
	the period 2015-2025, and we can see that it actually returns
	data about this temporal window. Notice how the query considers
	only the mean, which means that some metereological stations
	are excluded

	\item the Competency Question Q2 query can be found \href{https://github.com/zerise000/KG_Project2025_ClimateChange/blob/main/Evaluation/competencyQuestions/Q2.txt}{here}.
	The Competency Question asks for the distribution of the
	metereological facilities but since we cannot link directly
	the coordinates with a specific location we will just print
	the latitude and the longitude of the facility.
	
	\item The Competency Question Q3 query can be found \href{https://github.com/zerise000/KG_Project2025_ClimateChange/blob/main/Evaluation/competencyQuestions/Q3.txt}{here}.
.
	The body of the request splits the data in two distinct 
	periods, the one that goes from 2015 and terminates in
	the 2019, and the other one that goes from 2020 and terminates
	in 2025. After doing this it calculates the average temperature
	over these two periods separately and then it retrieves the
	temperature increase happened between the two eras. The expected 
	result should display the increase temperature over
	different areas in Trentino region, however since no areas
	are provided we consider the metereological stations as standalone
	locations

	\item The Competency Question Q5 query can be found \href{https://github.com/zerise000/KG_Project2025_ClimateChange/blob/main/Evaluation/competencyQuestions/Q5.txt}{here}.
.
	Also in this case we consider the metereological stations
	as distinct locations, which allows us to perform the query
	interrogation. The theoretical result should give us the
	different metereological stations, togheter with the statistics
	about the rainfalls, disposed in descendent order, and we can
	see that the actual query result actually satisfies the criteria 

	\item The Competency Question Q6 query can be found \href{https://github.com/zerise000/KG_Project2025_ClimateChange/blob/main/Evaluation/competencyQuestions/Q6.txt}{here}.
.
	The body of the request associates a number to each month
	in order to associate them with the season they belong, then
	it simply retrieves the interested data linked to a specific
	period. The expected result should return the months and the seasons over the
	considered decade with the highest solar radiation statistics,
	and we can see it actually satisfies the criteria

	\item The Competency Question Q8 query can be found \href{https://github.com/zerise000/KG_Project2025_ClimateChange/blob/main/Evaluation/competencyQuestions/Q8.txt}{here}.
.
	The expected result should return the humidity trend, 
	precisely the minimum,maximum and average values, over the
	summer of the last three years, which means that the query not
	only must consider only \textit{monthNum} equal to six, seven
	and eight (respectively June,July and August) but also it must filter
	the years not included in the period 2023-2025, extremes included.
	As we can see the query actually satisfies the required criteria

	\item The Competency Question Q9 query can be found \href{https://github.com/zerise000/KG_Project2025_ClimateChange/blob/main/Evaluation/competencyQuestions/Q9.txt}{here}.
. 
	We ask to the Knowledge Graph the rainfalls trend of the last
	five years by simply filtering the time when the data
	have been acquired and considering them only if they have
	relevations about the rain volumes, which means that the latter
	statistics in these cases must be greater than zero. The expected
	result should return the rainfall trend with the complete statistics,
	maximum, minimum and average rain volumes also in this case, and we
	can see how the query satisfies the required criteria

	\item The Competency Question Q10 query can be found \href{https://github.com/zerise000/KG_Project2025_ClimateChange/blob/main/Evaluation/competencyQuestions/Q10.txt}{here}..
	The query checks if the station contains the humidity measurements
	we are looking for, then it filters the data of the period
	2016-2025 (extremes included). Theoretically speaking it should
	return the humidity trends considering maximum, minimum and
	average values, togheter with the station that made the measurements
	and the year when the relevations have been produced, and we can see
	that the actual reult perfectly match the expectations
\end{itemize}
